\chapter{机器学习与 AI Alpha：从泛化到防泄漏 Pipeline}

\begin{itemize}
  \item \textbf{直接先修：}上册第 5、6、9、10、12、14、15、16、17 章；诊断入口见第 2.4 节。
  \item \textbf{学习产出：}交付一个时间边界可审计、具有简单基线、漂移响应和成本台账的固定推理 pipeline。
  \item \textbf{研究边界：}模型分数、概率校准和解释稳定性是三种不同证据；特征重要性不证明因果。
\end{itemize}

\section{经验风险、泛化与正则化}

给定损失 $\ell$ 和训练样本，经验风险为
\[
\widehat R(f)=\frac1n\sum_{i=1}^n\ell(y_i,f(x_i)),
\]
而泛化风险是对未来联合分布的期望 $R(f)=\mathbb E\ell(Y,f(X))$。训练误差下降并不保证二者差距缩小；模型选择、超参数搜索和预处理都会消耗研究选择权。

本章使用一个显式数据生成过程（DGP）
\[
Y=0.2X+0.8\mathbf 1\{X>0\}.
\]
目标由代码根据特征生成，不在 fixture 中预填模型预测；这使常数基线的误差可以独立复核。正则化目标
\[
\widehat R_\lambda(f)=\widehat R(f)+\lambda\Omega(f)
\]
把拟合与复杂度权衡写进同一目标。本章合成账本中训练风险为 $0.4256$，固定推理风险为 $1.5154$；取 $\lambda=0.01$、权重平方范数为 1，正则化目标为 $0.4356$。这些数只验证公式与数据边界。

训练、验证和固定推理必须沿时间顺序排列。随机切分把未来状态混入训练分布；全样本标准化则把未来均值与方差送入过去。两者都由实际索引或拟合截止日门禁拒绝。

\section{树模型与 boosting：从简单基线开始}

树模型用分裂规则形成分段常数函数；boosting 逐轮拟合当前损失的负梯度。选择平方损失、对数损失或排序损失之前，先说明目标究竟是收益、方向概率还是横截面次序。训练指标不能代替时间样本外指标。

代码在训练段实际拟合常数基线、单棵决策树桩和两轮残差 boosting。验证集 MSE 依次为 $1.0634$、$0.0650$ 和 $0.0425$，因此冻结 boosting。固定推理段的三个 MSE 依次为 $1.5154$、$0.2050$ 和 $0.1625$。这只是模型比较协议；如果为得到较小数值试了数百组参数，尝试次数和验证路径必须一起报告。

\begin{diagnostic}
只报告训练分数、只保留最佳超参数或在固定推理段继续 early stopping，都会使“样本外”名不副实。简单模型是基线，不是必须被复杂模型击败的装饰品。
\end{diagnostic}

\section{深度序列与表示模型}

序列输入约定为张量
\[
X\in\mathbb R^{B\times T\times F},
\]
分别表示 batch、时间步和特征。本章冻结形状为 $(2,3,2)$，mask 形状必须是 $(2,3)$，预测目标位于输入窗口之后一个时间步。padding mask、因果 attention mask 与缺失值 mask 不可混用。

非深度基线对最后一个有效时间步拟合线性模型；序列模型先按 mask 求表示，再经冻结随机投影、ReLU 和拟合输出层。二者只有在相同目标对齐、相同推理段和相同成本下比较才有意义。在冻结的合成任务上，两者 MSE 分别为 $2.866084$ 和 $0.403863$。这验证的是实际前向计算与标签对齐，并不证明序列模型在真实市场上占优。目标偏移为 0 会把窗口末值当作未来标签，门禁拒绝；未冻结随机种子则无法给出固定推理证据。

\begin{note}[Python 边界]
框架允许广播和自动 reshape，但“能运行”不说明 batch、time、feature 轴正确。每个张量边界都应保存形状、dtype、mask 语义和目标日期。
\end{note}

\section{漂移、校准与解释稳定性}

预测性能回答“损失多大”；概率校准回答“预测为 $p$ 的事件是否约有 $p$ 比例发生”；解释一致性回答“模型依赖的特征是否跨窗口稳定”。三者不能互相替代。

本章参考样本均值为 0，监控样本均值为 $0.8$，漂移距离超过阈值 $0.5$，触发停止自动上线并重新校准。二分类 Brier 分数
\[
\operatorname{BS}=\frac1n\sum_i(p_i-y_i)^2
\]
由 $0.16$ 降至 $0.025$，说明概率刻度改善，不证明收益增加。

训练与推理窗口的 Top-2 重要特征集合 Jaccard 只有 $1/3$。这要求报告解释不稳定；即使集合完全一致，特征重要性也只是模型依赖，不能宣称因果或结构稳定。

\section{固定推理与成本后结果}

两笔冻结交易毛收益为 $0.02$ 与 $0.01$，总成本 $0.006$，故
\[
R^{\mathrm{net}}=0.03-0.006=0.024.
\]
推理产物应绑定代码版本、fixture 哈希、模型配置、随机种子、预处理拟合截止日和目标日期。重新训练不能覆盖旧模型和旧预测。

\begin{implementationnote}
\begin{lstlisting}[language=bash]
uv run python notebooks/lower/ch03_ml_alpha.py \
  evidence/lower-ch03/oracle.json
\end{lstlisting}
\end{implementationnote}

\section{Route Capstone：防泄漏 Alpha pipeline}

Capstone 从干净环境复现训练、验证和固定推理，保留简单基线、超参数尝试账本、独立切分审计、漂移响应、校准报告、解释稳定性与成本台账。失败包必须可观察地复现随机切分、全样本预处理、未来目标对齐、随机性未冻结、只报告训练分数或超出调参预算，以及把特征重要性误写成因果，并展示修复后证据。

冻结摘要见 \path{reports/lower-ch03-summary.md}。限制报告必须明确数据制度、时间边界、模型漂移和交易摩擦，不能把特征重要性写成因果故事。

\section{四级问题}
\begin{enumerate}[leftmargin=*]
  \item \textbf{口述概念：}经验风险、泛化风险与正则化目标分别回答什么？
  \item \textbf{口述概念：}预测性能、概率校准和解释稳定性为何不能互相替代？
  \item \textbf{笔试推导：}计算本章三项风险与三个模型的 MSE。
  \item \textbf{笔试推导：}写出 Brier 分数并核对校准前后结果。
  \item \textbf{数值编程：}扩展张量到四步窗口，构造正确 mask 与下一期目标。
  \item \textbf{数值编程：}加入第三笔交易并独立重算成本后收益。
  \item \textbf{研究判断：}为什么随机切分、全样本 scaler 和推理段 early stopping 都是泄漏？
  \item \textbf{研究判断：}漂移触发后应停止、重校准还是重训？写出可审计协议。
\end{enumerate}

